\section{Introduction}
\label{sec:introduction}

Flow matching (FM) trains continuous-time generative models by regressing a neural velocity field
against conditional vector fields along prescribed probability paths
\citep{lipman2023flow,tong2023conditional,albergo2023stochastic}. After training, sampling
requires solving an ordinary differential equation (ODE). If the trained velocity is
$v_\theta:\mathbb{R}^d \times [0,1] \to \mathbb{R}^d$, generation from a source sample
$x_0 \sim p_0$ is obtained by integrating
\begin{equation}
  \frac{d x(t)}{d t} = v_\theta(x(t),t), \qquad x(0)=x_0.
  \label{eq:fm-ode}
\end{equation}
The computational cost of sampling is therefore dominated by neural function evaluations (NFE).
Low-NFE generation is attractive, but it exposes a basic numerical question: where should a small
number of time steps be placed?

Several lines of work address this question indirectly. Rectified flows and optimal-transport
conditional flow matching attempt to learn straighter or better coupled paths
\citep{liu2022rectified,tong2023conditional}. Acceleration-regularized training methods, including
recent Iso-FM-style objectives, penalize the material acceleration of the learned field during
training \citep{khan2026isofm}. Learned and bespoke solvers optimize the update rule or solver
parameters for a trained generative ODE \citep{shaul2024bespoke,shaul2024bns}. In diffusion
sampling, trajectory-regularity and adaptive schedule methods choose nonuniform timesteps from
properties of probability-flow trajectories \citep{chen2024trajectory,chen2024adaptive}. These
works make clear that nonuniform timesteps are not new. The question we study is narrower:
can a standard flow-matching checkpoint supply a cheap, native, training-free diagnostic for where
Euler steps should be spent?

This paper studies two parallel population-calibrated meshes. Self-Curvature Time Warping (SCTW)
performs a fine, low-batch probe rollout, estimates how fast the velocity changes along generated
trajectories, and places sampling knots at quantiles of a tempered curvature density. The
diagnostic is the material derivative of the learned velocity field,
\begin{equation}
  \frac{D v_\theta}{D t}(x,t)
  =
  \partial_t v_\theta(x,t)
  +
  \nabla_x v_\theta(x,t)\, v_\theta(x,t),
  \label{eq:material-derivative}
\end{equation}
but the implemented SCTW estimator avoids Jacobian--vector products: it measures finite
differences of the velocity along the model's own rollout. Adjoint-SCTW asks a second question:
how strongly will a local Euler defect affect the terminal sample? It propagates terminal
sensitivity backward through the same fine pilot with randomized state adjoints and weights the
local profile accordingly. The checkpoint remains frozen, and one population mesh is reused for
all subsequent samples.

The contribution is intentionally modest and, we believe, scientifically useful:
\begin{itemize}
  \item We derive a time-mesh rule for Euler sampling from the local truncation error induced by
  the material derivative of a learned FM velocity.
  \item We formulate terminally transported Euler defects through the flow-map tangent and derive
  finite-difference, exact JVP, and randomized discrete-adjoint estimators of the same profile.
  \item We connect the local material acceleration to an exact pressure/momentum identity for linear-interpolant
  FM, explaining why marginal trajectories can curve even when conditional paths are straight.
  \item We compare both adaptive meshes with uniform Euler, matched-NFE Heun, VSFM schedules,
  early-, late-, and symmetric-heavy meshes, and a low-dimensional offline schedule optimizer.
  \item We separate learned-ODE integration error from endpoint distribution quality. Their strong
  alignment on staged-shapes and weak or negative alignment on CIFAR explains both the principal
  positive result and the principal limitation.
\end{itemize}
